\chapter*{Preface: What happens if you flip the exponents?}
\addcontentsline{toc}{chapter}{Preface}
\markboth{PREFACE}{PREFACE}

This book was not born at a desk. It began in a child's bedroom, in
the dark, while I was putting my daughter to sleep. I lay
next to her, determined not to doze off myself -- only to keep her
company -- and nothing keeps a mind awake as reliably as a thought
that won't let go. For me such thoughts are numbers; my head
computes on its own. That evening I was practising mental
arithmetic, out of a slightly stubborn resolve not to unlearn
calculation in the age of computers: 728 divided by 4. The 700 is 7
times 100, hence 7 times 4 times 25 -- divided by four that leaves
7 times 25, which is 175; the 28 trivially gives 7; together, 182.

Nothing about this is remarkable, except for what caught my
attention while taking the number apart: how naturally I sort
hundreds, tens, and units into separate boxes. From there it was a
short walk through everything I knew about positional systems --
binary numbers, the base-twelve system you can count on the
segments of your fingers -- to a question I had never seen asked
anywhere: what happens if the counting bases change \emph{inside} a
single number? That such systems already existed -- the factorial
numbers -- I did not know that night; the first literature search
turned them up later. But ascending bases struck me as boring. I
wanted it the other way round: the most possibilities at the far
right.\footnote{Strictly speaking no exponent gets flipped; what is
reversed is the direction in which the bases grow -- the weight
ordering of the positional system. The title follows the original
thought, not the terminology.} Whoever takes that thought seriously
runs straight into a necessity: such a system cannot count to
infinity -- infinity would have been nonsense here, so: finite.
What at first looks like a defect is the birth of this book's
central notion: every number lives inside a \emph{horizon}, a
finite space with a fixed number of places. My daughter was long
asleep by then. My wife told me later I had been upstairs for a
good hour;
it had felt like twenty-five minutes.

That leading zeros change the numbers in such a system -- the
oddity that drives everything that follows -- I did not yet know
that night. It came to light the way ideas get tested nowadays: I
sat down at the computer and explained the thought to GPT, asking
it to check whether all of this already existed. The first verdict
was dry: nonsense. But we kept at it, and after a while we realized
together that it was not complete nonsense after all. In those
conversations the zeros surfaced -- zeros that suddenly mean
something at the front of a number -- along with many of the ideas
this book picks up. One sentence from that time has stayed with me:
4 times 5 can be an area -- and if the 4 is a scalar, a length; in
both cases 20. It has only a marginal bearing on Horimetrik, but it
made an impression -- and it finds a late echo in
Chapter~\ref{ch:strukturpolynome}. The name, by the way, was coined
entirely without ceremony: ``horizon arithmetic''
(\emph{Horizont-Arithmetik}) was too long for me. So: Horimetrik.

At some point the subject had outgrown casual conversation, and GPT
and I agreed that it deserved an agent of its own. That is how
Fable entered the race -- the Claude model by Anthropic that
appears on the title page as mathematical collaborator. Let me be
honest: mathematically, I was left behind fast. My computer science
degree taught me mathematics, minor subject included -- but that
was more than a decade ago. I understand the outlines, and the
outlines were exciting enough: not because there was any
application, but because it all seemed novel enough to be worth
pursuing. Many fields of mathematics began without an application.
Fable can do mathematics; so I had everything recomputed and a
clean theory built, and over time more and more bridges to
established mathematics appeared. GPT kept contributing ideas along
the way. My thanks go to both -- read for yourselves what the AI
and I have found.

The book pursues two goals. First, it wants to develop the
mathematics of this space honestly -- with definitions, proofs, and
also with the wrong turns we took along the way. Second, it wants
to stay readable: every technical passage returns to the surface
before the next one begins. Whoever reads only the formula-free
sections should still understand what is going on.

Whether Horimetrik becomes a theory in its own right or remains a
particularly instructive plaything was something we ourselves did
not know while writing. That is precisely what this book is about.
